\documentclass[11pt]{article}
\usepackage[T1]{fontenc}
\usepackage{textcomp}
\usepackage[margin=0.75in]{geometry}
\usepackage{amsmath,amssymb,amsthm}
\usepackage{array,booktabs}
\usepackage{hyperref}
\setlength{\parindent}{0pt}
\newtheorem{theorem}{Theorem}
\theoremstyle{definition}
\newtheorem{definition}{Definition}
\title{Flow Proof Artifact: Subgroup}
\author{Generated by \texttt{flow doc proof}}
\date{Flow Proof Book}
\begin{document}
\maketitle
Subgroup axioms as a subset closed under multiplication and inverses.\\[0.5em]
\textbf{Source.} dummit-foote — *Abstract Algebra*, §2.2\\[0.5em]
\bigskip
\noindent{\large \textbf{Definition 1.} \textit{A subgroup contains the group identity} \hfill Subgroup \cdot membership \cdot the identity lies in every subgroup}\par
\smallskip\noindent\textit{Coordinate: Subgroup \cdot membership \cdot the identity lies in every subgroup}\par
\smallskip\noindent\textit{Source: dummit-foote}\par
\medskip
\noindent\textbf{Goal.} A subgroup contains the group identity\par
\noindent$\displaystyle \forall H \in Subgroup\quad 1 in H$\par
\medskip
\noindent
\begin{tabular}{@{} >{\raggedright\arraybackslash}p{0.06\textwidth} >{\raggedright\arraybackslash}p{0.47\textwidth} >{\raggedleft\arraybackslash}p{0.41\textwidth} @{}}
\textbf{\#} & \textbf{Proof} & \textbf{Mathematics} \\
\midrule
\textbf{1.} & We stipulate the identity lies in every subgroup for membership on Subgroup — this is a definition, not a derived fact. & \\[0.45em]
\textbf{2.} & This follows directly from the definition: 1 in H. Hence proven. & $\displaystyle 1 in H$ \\[0.45em]
\bottomrule
\end{tabular}
\label{thm:Subgroup-membership-the_identity_lies_in_every_subgroup}
\par\medskip\hrule
\bigskip
\noindent{\large \textbf{Definition 2.} \textit{A subgroup is closed under multiplication} \hfill Subgroup \cdot multiplication \cdot closed under multiplication}\par
\smallskip\noindent\textit{Coordinate: Subgroup \cdot multiplication \cdot closed under multiplication}\par
\smallskip\noindent\textit{Source: dummit-foote}\par
\medskip
\noindent\textbf{Goal.} A subgroup is closed under multiplication\par
\noindent$\displaystyle \forall H \in Subgroup \forall a \in Group \forall b \in Group\quad a \cdot b in H$\par
\medskip
\noindent
\begin{tabular}{@{} >{\raggedright\arraybackslash}p{0.06\textwidth} >{\raggedright\arraybackslash}p{0.47\textwidth} >{\raggedleft\arraybackslash}p{0.41\textwidth} @{}}
\textbf{\#} & \textbf{Proof} & \textbf{Mathematics} \\
\midrule
\textbf{1.} & We stipulate closed under multiplication for multiplication on Subgroup — this is a definition, not a derived fact. & \\[0.45em]
\textbf{2.} & We split into exhaustive cases — the claim must hold in each one. & \\[0.45em]
\textbf{3.} & Case 1 (see step 2): suppose a in H. & \\[0.45em]
\textbf{4.} & Case 2 (see step 2): suppose b in H. & \\[0.45em]
\textbf{5.} & From step 4, this implies a times b in H. Together with the other cases (step 3 and step 4), the goal is discharged. Hence proven. & $\displaystyle a \cdot b in H$ \\[0.45em]
\bottomrule
\end{tabular}
\label{thm:Subgroup-multiplication-closed_under_multiplication}
\par\medskip\hrule
\bigskip
\noindent{\large \textbf{Definition 3.} \textit{A subgroup is closed under taking inverses} \hfill Subgroup \cdot inverse \cdot closed under inverses}\par
\smallskip\noindent\textit{Coordinate: Subgroup \cdot inverse \cdot closed under inverses}\par
\smallskip\noindent\textit{Source: dummit-foote}\par
\medskip
\noindent\textbf{Goal.} A subgroup is closed under taking inverses\par
\noindent$\displaystyle \forall H \in Subgroup \forall g \in Group\quad inv(g) in H$\par
\medskip
\noindent
\begin{tabular}{@{} >{\raggedright\arraybackslash}p{0.06\textwidth} >{\raggedright\arraybackslash}p{0.47\textwidth} >{\raggedleft\arraybackslash}p{0.41\textwidth} @{}}
\textbf{\#} & \textbf{Proof} & \textbf{Mathematics} \\
\midrule
\textbf{1.} & We stipulate closed under inverses for inverse on Subgroup — this is a definition, not a derived fact. & \\[0.45em]
\textbf{2.} & We split into exhaustive cases — the claim must hold in each one. & \\[0.45em]
\textbf{3.} & Case 1 (see step 2): suppose g in H. & \\[0.45em]
\textbf{4.} & From step 3, this implies inv(g) in H. Together with the other cases (step 3), the goal is discharged. Hence proven. & $\displaystyle inv(g) in H$ \\[0.45em]
\bottomrule
\end{tabular}
\label{thm:Subgroup-inverse-closed_under_inverses}
\par\medskip\hrule
\end{document}
